% GENERATED FILE. Edit canonical lessons, then run npm run generate:books.
% canonical-source-hash: fnv1a64:ffeac82100ed3594
% canonical-lessons: KA-C464-upahara, KA-C464-menu, KA-C464-paricaraka, KA-C464-bhaksisu, KA-C464-ardar-madu

\chapter{Breakfast at the Hotel}
\label{ch:ka-breakfast-at-the-hotel}
\hlchaptermodality{464}

\begin{chapteropening}
\textbf{By the end of this chapter:} \emph{I can say breakfast, a menu, a waiter, a gratuity and to order.}

Complete the last lesson of chapter 464: I can say breakfast, a menu, a waiter, a gratuity and to order.
\end{chapteropening}

\section[\texorpdfstring{\kn{ಉಪಹಾರ} (upahāra)}{ಉಪಹಾರ (upahāra)}]{\kn{ಉಪಹಾರ} (upahāra) — breakfast; a light meal}

\label{lesson:KA-C464-upahara}



\begin{quote}
Before the new one: say the Kannada for a water heater, then the Kannada for a resort.
\end{quote}

\subsection*{You'll want to know: \kn{ಉಪಹಾರ}}
\textbf{\kn{ಉಪಹಾರ}} — \emph{upahāra} — \textquotedblleft{}breakfast; a light meal\textquotedblright{}.

\begin{quote}
\kn{ಹೋಟೆಲ್ನಲ್ಲಿ} \kn{ಉಪಹಾರಕ್ಕೆ} \kn{ಇಡ್ಲಿ} \kn{ಮತ್ತು} \kn{ದೋಸೆ} \kn{ಇತ್ತು}.
\end{quote}

\begin{quote}
\emph{hōṭelnalli upahārakke iḍli mattu dōse ittu.}
\end{quote}

\begin{quote}
At the hotel there were idlis and dosas for breakfast.
\end{quote}

\kn{ಬೆಳಗಿನ} \kn{ಉಪಹಾರ}, morning breakfast; \kn{ತಿಂಡಿ} is the everyday word.

\begin{cousinweb}
\textbf{\kn{ಉಪಹಾರ}} (\emph{upahāra}) is borrowed from Sanskrit: Monier-Williams gives \emph{upahāra} as \textquotedblleft{}an offering, a present\textquotedblright{}. The sense of a light meal is a later Indian one.
\end{cousinweb}

\begin{grammarlens}[title={What you've built}]
One more word for getting around.
\end{grammarlens}

\subsection*{Your turn}
\begin{itemize}
\raggedright
  \item \emph{Say it:} \emph{upahāra}
  \item \emph{Say it:} \emph{upahāra}, once more
  \item \emph{Say it:} \emph{hōṭelnalli upahārakke iḍli mattu dōse ittu.} — At the hotel there were idlis and dosas for breakfast.
  \item \emph{Recall:} say the Kannada for a water heater, then the Kannada for a resort, then say \emph{upahāra} again
\end{itemize}

\begin{tcolorbox}[breakable,colback=teal!4,colframe=teal!35!black,title={Before you move on}]
What does \kn{ಉಪಹಾರ} mean? (\textquotedblleft{}Breakfast; a light meal\textquotedblright{}.) Say it once more.
\end{tcolorbox}
\section[\texorpdfstring{\kn{ಮೆನು} (menu)}{ಮೆನು (menu)}]{\kn{ಮೆನು} (menu) — a menu}

\label{lesson:KA-C464-menu}



\begin{quote}
Before the new one: say the Kannada for a resort, then the Kannada for breakfast.
\end{quote}

\subsection*{You'll want to know: \kn{ಮೆನು}}
\textbf{\kn{ಮೆನು}} — \emph{menu} — \textquotedblleft{}a menu\textquotedblright{}.

\begin{quote}
\kn{ಮೆನುವಿನಲ್ಲಿ} \kn{ತುಂಬಾ} \kn{ದೋಸೆಗಳಿವೆ}.
\end{quote}

\begin{quote}
\emph{menuvinalli tumbā dōsegaḷive.}
\end{quote}

\begin{quote}
There are lots of dosas on the menu.
\end{quote}

The ending -\kn{ವಿನಲ್ಲಿ}: \kn{ಮೆನು} takes -\kn{ವ}- before a case ending.

\begin{cousinweb}
\textbf{\kn{ಮೆನು}} (\emph{menu}) is a loan from English \emph{menu}.
\end{cousinweb}

\begin{grammarlens}[title={What you've built}]
One more word for getting around.
\end{grammarlens}

\subsection*{Your turn}
\begin{itemize}
\raggedright
  \item \emph{Say it:} \emph{menu}
  \item \emph{Say it:} \emph{menu}, once more
  \item \emph{Say it:} \emph{menuvinalli tumbā dōsegaḷive.} — There are lots of dosas on the menu.
  \item \emph{Recall:} say the Kannada for a resort, then the Kannada for breakfast, then say \emph{menu} again
\end{itemize}

\begin{tcolorbox}[breakable,colback=teal!4,colframe=teal!35!black,title={Before you move on}]
What does \kn{ಮೆನು} mean? (\textquotedblleft{}A menu\textquotedblright{}.) Say it once more.
\end{tcolorbox}
\section[\texorpdfstring{\kn{ಪರಿಚಾರಕ} (paricāraka)}{ಪರಿಚಾರಕ (paricāraka)}]{\kn{ಪರಿಚಾರಕ} (paricāraka) — a waiter; an attendant}

\label{lesson:KA-C464-paricaraka}



\begin{quote}
Before the new one: say the Kannada for breakfast, then the Kannada for a menu.
\end{quote}

\subsection*{You'll want to know: \kn{ಪರಿಚಾರಕ}}
\textbf{\kn{ಪರಿಚಾರಕ}} — \emph{paricāraka} — \textquotedblleft{}a waiter; an attendant\textquotedblright{}.

\begin{quote}
\kn{ಪರಿಚಾರಕ} \kn{ನಮಗೆ} \kn{ಬಿಸಿ} \kn{ಕಾಫಿ} \kn{ತಂದನು}.
\end{quote}

\begin{quote}
\emph{paricāraka namage bisi kāphi tandanu.}
\end{quote}

\begin{quote}
The waiter brought us hot coffee.
\end{quote}

Plural: \kn{ಪರಿಚಾರಕರು}, \emph{paricārakaru}.

A formal word; people often call out \kn{ಸರ್} or \kn{ಅಣ್ಣ} to a waiter.

\begin{cousinweb}
\textbf{\kn{ಪರಿಚಾರಕ}} (\emph{paricāraka}) is borrowed from Sanskrit \emph{paricāraka}, \textquotedblleft{}an assistant, an attendant\textquotedblright{} (Monier-Williams), from \emph{pari-car}, \textquotedblleft{}to move about, to attend on\textquotedblright{}.
\end{cousinweb}

\begin{grammarlens}[title={What you've built}]
One more word for getting around.
\end{grammarlens}

\subsection*{Your turn}
\begin{itemize}
\raggedright
  \item \emph{Say it:} \emph{paricāraka}
  \item \emph{Say it:} \emph{paricāraka}, once more
  \item \emph{Say it:} \emph{paricāraka namage bisi kāphi tandanu.} — The waiter brought us hot coffee.
  \item \emph{Recall:} say the Kannada for breakfast, then the Kannada for a menu, then say \emph{paricāraka} again
\end{itemize}

\begin{tcolorbox}[breakable,colback=teal!4,colframe=teal!35!black,title={Before you move on}]
What does \kn{ಪರಿಚಾರಕ} mean? (\textquotedblleft{}A waiter; an attendant\textquotedblright{}.) Say it once more.
\end{tcolorbox}
\section[\texorpdfstring{\kn{ಭಕ್ಷೀಸು} (bhakṣīsu)}{ಭಕ್ಷೀಸು (bhakṣīsu)}]{\kn{ಭಕ್ಷೀಸು} (bhakṣīsu) — a gratuity; a tip}

\label{lesson:KA-C464-bhaksisu}



\begin{quote}
Before the new one: say the Kannada for a menu, then the Kannada for a waiter.
\end{quote}

\subsection*{You'll want to know: \kn{ಭಕ್ಷೀಸು}}
\textbf{\kn{ಭಕ್ಷೀಸು}} — \emph{bhakṣīsu} — \textquotedblleft{}a gratuity; a tip\textquotedblright{}.

\begin{quote}
\kn{ನಾವು} \kn{ಪರಿಚಾರಕನಿಗೆ} \kn{ಹತ್ತು} \kn{ರೂಪಾಯಿ} \kn{ಭಕ್ಷೀಸು} \kn{ಕೊಟ್ಟೆವು}.
\end{quote}

\begin{quote}
\emph{nāvu paricārakanige hattu rūpāyi bhakṣīsu koṭṭevu.}
\end{quote}

\begin{quote}
We gave the waiter a ten rupee tip.
\end{quote}

\kn{ಭಕ್ಷೀಸು} \kn{ಕೊಡು}, to tip.

\begin{cousinweb}
\textbf{\kn{ಭಕ್ಷೀಸು}} (\emph{bhakṣīsu}) is usually traced to Persian \emph{bakhshish}, \textquotedblleft{}a gift, a present\textquotedblright{}, which reached Kannada through Hindustani.
\end{cousinweb}

\begin{grammarlens}[title={What you've built}]
One more word for getting around.
\end{grammarlens}

\subsection*{Your turn}
\begin{itemize}
\raggedright
  \item \emph{Say it:} \emph{bhakṣīsu}
  \item \emph{Say it:} \emph{bhakṣīsu}, once more
  \item \emph{Say it:} \emph{nāvu paricārakanige hattu rūpāyi bhakṣīsu koṭṭevu.} — We gave the waiter a ten rupee tip.
  \item \emph{Recall:} say the Kannada for a menu, then the Kannada for a waiter, then say \emph{bhakṣīsu} again
\end{itemize}

\begin{tcolorbox}[breakable,colback=teal!4,colframe=teal!35!black,title={Before you move on}]
What does \kn{ಭಕ್ಷೀಸು} mean? (\textquotedblleft{}A gratuity; a tip\textquotedblright{}.) Say it once more.
\end{tcolorbox}
\section[\texorpdfstring{\kn{ಆರ್ಡರ್} \kn{ಮಾಡು} (ārḍar māḍu)}{ಆರ್ಡರ್ ಮಾಡು (ārḍar māḍu)}]{\kn{ಆರ್ಡರ್} \kn{ಮಾಡು} (ārḍar māḍu) — to order}

\label{lesson:KA-C464-ardar-madu}



\begin{quote}
Before the new one: say the Kannada for a waiter, then the Kannada for a gratuity.
\end{quote}

\subsection*{You'll want to know: \kn{ಆರ್ಡರ್} \kn{ಮಾಡು}}
\textbf{\kn{ಆರ್ಡರ್} \kn{ಮಾಡು}} — \emph{ārḍar māḍu} — \textquotedblleft{}to order\textquotedblright{}.

\begin{quote}
\kn{ನಾನು} \kn{ಎರಡು} \kn{ಕಾಫಿ} \kn{ಆರ್ಡರ್} \kn{ಮಾಡಿದೆ}.
\end{quote}

\begin{quote}
\emph{nānu eraḍu kāphi ārḍar māḍide.}
\end{quote}

\begin{quote}
I ordered two coffees.
\end{quote}

Past (I …): \kn{ಆರ್ಡರ್} \kn{ಮಾಡಿದೆ}, \emph{ārḍar māḍide}.

Food in a hotel, or goods online: \kn{ಆನ್ಲೈನ್ನಲ್ಲಿ} \kn{ಆರ್ಡರ್} \kn{ಮಾಡು}.

\begin{cousinweb}
\textbf{\kn{ಆರ್ಡರ್} \kn{ಮಾಡು}} (\emph{ārḍar māḍu}) joins a loan from English \emph{order} and the native verb \kn{ಮಾಡು}, to do.
\end{cousinweb}

\begin{grammarlens}[title={What you've built}]
One more word for getting around.
\end{grammarlens}

\subsection*{Your turn}
\begin{itemize}
\raggedright
  \item \emph{Say it:} \emph{ārḍar māḍu}
  \item \emph{Say it:} \emph{ārḍar māḍu}, once more
  \item \emph{Say it:} \emph{nānu eraḍu kāphi ārḍar māḍide.} — I ordered two coffees.
  \item \emph{Recall:} say the Kannada for a waiter, then the Kannada for a gratuity, then say \emph{ārḍar māḍu} again
\end{itemize}

\begin{tcolorbox}[breakable,colback=teal!4,colframe=teal!35!black,title={Before you move on}]
What does \kn{ಆರ್ಡರ್} \kn{ಮಾಡು} mean? (\textquotedblleft{}To order\textquotedblright{}.) Say it once more.
\end{tcolorbox}
